\textbf{Launch together.} After movement time let children move two counters freely for a minute. On a small non-task graph, put C down first, let a child place R, and show one-road C motion. Flip the turn tile. R stays; flip the tile again. Both pieces remain visible throughout. Show that sharing a dot ends the game immediately. Ask the pair to try catching or escaping, then distribute the first selected board. Do not explain the tree/cycle strategies before play.

\textbf{K--1 route (adult-led).}
Run the same launch at the K--1 table with one pair kit and shared pp.~1 and 3. Play the five-dot path P5 aloud: the parent asks ``Where should C start so that C always catches R?'', turns the tile after every move, and lets the children place and move both counters. Then set the triangle R3 beside the square R4 and ask ``On which ring can R get away?'' Plain-words key: on the path, C starts in the middle dot c and catches R within two cop moves wherever R goes. On the triangle C catches R on the very first move, from any dot. On the square R gets away forever by standing across from C and stepping to the new across-dot whenever C moves. The adult reads and keeps the tile; the children choose every move. Whether K--1 children keep stays and turn order is a prediction, not an observation; if the turns break down, stay with the path.
